\documentclass[10.5pt,a4paper]{article}
\usepackage[margin=0.75in,top=0.6in,bottom=0.6in]{geometry}
\usepackage[T1]{fontenc}
\usepackage[utf8]{inputenc}
\usepackage{setspace}
\usepackage{xcolor}
\usepackage{amsmath}
\usepackage{graphicx}
\usepackage{fancyhdr}

\definecolor{headerblue}{RGB}{30,80,140}
\definecolor{midgray}{RGB}{180,180,180}
\definecolor{inputcolor}{RGB}{235,247,235}
\setstretch{0.97}

\pagestyle{fancy}
\fancyhf{}
\fancyhead[L]{\color{headerblue}\small\textbf{IITISoC 2025}\ |\ Mental-Health Digital-Twin Project}
\fancyhead[R]{\color{midgray}\small June 2026}
\renewcommand{\headrulewidth}{0.4pt}
\renewcommand{\headrulecolor}{\color{headerblue}}
\fancyfoot[C]{\color{midgray}\thepage}

\newcommand{\stagesep}{{\color{headerblue}\hrule width \linewidth height 1.2pt}\vspace{0.4em}}
\newcommand{\stagetitle}[2]{%
  \vspace{0.4em}
  {\color{headerblue}\Huge\bfseries #1}\\[0.15em]
  {\large\color{midgray}#2}\\[0.3em]
  {\color{headerblue}\hrule width \linewidth height 0.6pt}\vspace{0.6em}}

\begin{document}

\stagesep
\stagetitle{Dashboard Outputs --- Stage 4, Stage 5, and XAI}{}

The live dashboard is the presentation layer of the mental-health digital twin: a React clinical front-end, served through Vercel, that reaches the Flask inference backend over a Cloudflare tunnel. Every widget is wired to a real pipeline output rather than a static visualisation, so the numbers a clinician sees on screen are the same tensors produced by the anomaly ensemble, the calibrated risk model, and the explainability engine described in the earlier stages.

\section*{Stage 4 --- Anomaly Detection Outputs}

The centrepiece of the clinical view is the behavioural risk gauge, a 0--100 score derived from the weighted consensus of the four Stage 4 detectors --- Mahalanobis distance, Gaussian copula, isolation forest, and KNN distance --- fused with the adaptive ensemble weights. As the fused anomaly probability $p_{\text{anom}}$ crosses the adaptive thresholds derived from the 97th percentile of each user's history, the gauge changes both colour and label, cycling from \emph{Nominal baseline} through \emph{Mild elevation} and \emph{Moderate anomaly} to \emph{Critical threshold violation}, so a sustained drift is visible at a glance.

Directly beneath it, the \emph{What's driving that signal} panel renders the four detectors as synchronised sparklines, each drawn in a colour that moves from cyan through amber to orange as the per-detector score climbs past the 30\% and 60\% bands. Hovering anywhere on a sparkline snaps a crosshair to the nearest day and prints the exact score with its semantic label (Stable, Moderate, or Elevated), letting the clinician attribute a rising risk to its mechanistic cause. Beside these, the CUSUM panel plots the upper and lower cumulative sums $S_t^{+}$ and $S_t^{-}$ against the decision threshold $h$, and the moment either accumulation crosses $h$ the chart flags a persistent baseline drift; both panels support zooming and scrolling so any time window can be inspected. Whenever any detector exceeds 40\%, the platform also pushes a typed alert (info, warning, or critical) into the notification bell, giving an at-a-glance audit trail of when and why risk began to climb.

\section*{Stage 5 --- Calibrated Risk Outputs}

The Stage 5 output is presented as a temperature-calibrated probability $p_{\text{cal}}$, mapped back onto the 0--100 risk scale and paired with the raw, uncalibrated probability so that the effect of the temperature scaling learned on the held-out cohort is always visible side by side. The card simultaneously reports the risk level (LOW, MODERATE, or HIGH), the calibrated model confidence, and the intervention flag, which activates whenever the calibrated probability crosses the deployment threshold and recommends clinical follow-up within 48 hours. A set of natural-language insight bullets then binds the numbers to the user's own data, quoting, for example, the sleep-duration disruption percentage and the number of journal entries processed by the pipeline.

\section*{Explainable AI Outputs (TreeSHAP)}

The Explainable AI tab renders the TreeSHAP explanation produced over the 2,336-dimensional feature vector. The explanation is anchored at the model's base value, and the top contributing features are drawn as horizontal bars whose width is proportional to the absolute Shapley value $|\phi_j|$; features that push risk upward are coloured red while protective features are drawn green, with the signed value $\phi_j$ printed to four decimal places. These individual attributions are then aggregated into domain-level impacts, shown as stacked percentages, so the clinician can see, for instance, how much of the total SHAP contribution comes from sleep-hygiene features versus social-output features. Finally, a set of natural-language sentences translates the attributions into clinical statements, and the whole panel is labelled with the transparency score and the inference latency of the explanation computation.

\section*{Patient Portal --- Daily Input and Feedback}

On the patient side, the Daily Alignment Portal turns raw lived experience into the feature stream that feeds the later stages. Each day the user submits a journal file (TXT, MD, DOCX, or PDF), optionally a voice recording, or an exported WhatsApp, Telegram, or Reddit conversation, which is parsed, date-matched to the user's own messages, and previewed before import. Four sliders --- sleep hours, sleep quality, activity level, and mood --- capture the structured lifestyle signal. A calibration progress bar tracks the entry count toward the minimum of 14, and once the baseline is complete the portal displays a \emph{Baseline Calibrated} banner confirming that the personal model is active. The entry-history table closes the loop, showing which submissions were feature-extracted so the clinician can verify the data quality behind every risk score rendered on the dashboard.

\end{document}
